\section{Animación}
Con el propósito de evaluar la viabilidad de la extensión de Blender como herramienta para la generación 
de material didáctico, se diseñó y produjo una animación orientada a la divulgación técnica.
\medskip

La pieza audiovisual está dirigida a una audiencia familiarizada con los fundamentos
 del aprendizaje automático y las redes neuronales densas, teniendo como objetivo principal 
 introducir conceptos nucleares de las redes neuronales convolucionales, tales como:
\begin{itemize}
    \item La operación formal de convolución bidimensional.
    \item La arquitectura jerárquica de la red (capas convolucionales y de reducción espacial o \textit{MaxPool2d}).
    \item Las divergencias estructurales respecto a redes densamente conectadas.
    \item El campo receptivo y la evolución representacional entre capas tempranas y profundas.
    \item Las diferencias funcionales entre tareas de clasificación y detección de objetos.
\end{itemize}

Dado su carácter formativo, la animación adopta un lenguaje visual depurado y sobrio para priorizar la
comprensión de las estructuras de datos tridimensionales. Para su realización se articularon los componentes 
procedimentales desarrollados en este trabajo junto con herramientas externas de posproducción digital.
\medskip

Los detalles más específicos sobre el contenido de la animación se encuentran en el Anexo~\ref{anx:guion_anim}.
% Agregar discusión técnica de la animación cuando quede
